\section{Método}
\subsection{Desenho e escopo}
Foi realizada análise documental piloto de caso único com inspeção estática de código, conforme o protocolo versionado do projeto \rep. A escolha do caso decorre da disponibilidade do GDD e do repositório; as seis unidades foram selecionadas intencionalmente para exercitar correspondência explícita, divergência literal e falta de parâmetro, sem procedimento probabilístico ou pretensão de representatividade do documento \rep.

O piloto permite verificar se as categorias produzem registros compreensíveis e localizáveis, mas não estima a frequência de problemas no jogo inteiro \rep. Não houve questionário, entrevista, observação sistemática de estudantes, comparação de aprendizagem ou coleta de notas no corpus descrito \rep. A eventual ampliação para essas dimensões configurará outro escopo e precisará de protocolo próprio, em lugar de inferência retrospectiva a partir deste material (Brasil, 2016; Night Heist, 2026).

\subsection{Corpus e controle de versões}
O corpus primário reúne a leitura textual do PDF fornecido pelo usuário e os arquivos do baseline identificado pelo commit \nolinkurl{72699f0c2752d029d8e8c3b48420c18a70dbe2a1} \fonte. A identificação por commit fixa o objeto técnico e evita que mudanças posteriores sejam incorporadas silenciosamente à análise \rep. Os instrumentos científicos, por sua vez, são documentos de trabalho da branch \texttt{docs/projeto-academico-cientifico}, proposta na PR 2, e não arquivos integrantes desse baseline do jogo \rep.

O arquivo no Drive é denominado \emph{GDD -- Night Heist Escape Phase (Versão 0).pdf}, enquanto o conteúdo identifica \emph{Night Heist: Escape Phase V2}, versão 0.2, MVP sem grade descendente; as duas identificações foram preservadas no registro de fonte \gdd. Os metadados consultados informaram tamanho de 142006 bytes e modificação em 8 de outubro de 2026, mas os bytes originais não foram materializados no ambiente de análise \rep.

Consequentemente, não foi calculado hash do PDF nem validada sua paginação, e a inspeção visual do mockup permanece pendente \rep. Os localizadores utilizados são seções e identificadores de mecânica, adequados à leitura textual disponível, sem inventar números de página \gdd. O artigo incorpora o GDD como fonte e objeto de análise, não como reprodução integral anexada \rep.

A documentação interna de design, backlog e validação foi tratada como fonte auxiliar do repositório, sem substituir o GDD original \rep. Conversas com o agente e registros que individualizem decisões não compõem o corpus atual, motivo pelo qual autoria de parâmetros não foi deduzida a partir do código ou do relato geral da experiência \rep.

\subsection{Instrumento e categorias}
A matriz registra identificador da unidade, origem, localizador e trecho do GDD, critério, localizador da implementação, teste associado quando existente, evidência, categoria, ambiguidade, responsável, fonte da decisão, estado técnico e revisão humana \rep. A presença de coluna de teste não significa execução: nos seis registros iniciais, o estado é inspeção estática e a revisão humana está pendente \rep.

\begin{singlespace}
\begin{longtable}{p{0.22\textwidth}p{0.69\textwidth}}
\caption{Categorias operacionais de correspondência.}\label{tab:categories}\\
\toprule Categoria & Regra adotada no instrumento \\\midrule\endfirsthead
\toprule Categoria & Regra adotada no instrumento \\\midrule\endhead
Preservado & Todos os critérios explícitos da unidade são atendidos na evidência examinada.\\
Parcial & Parte dos critérios é atendida; as partes não atendidas precisam ser identificadas.\\
Alterado & Há implementação correspondente que contradiz pelo menos um critério explícito.\\
Omitido & Unidade explícita aplicável não tem implementação após busca e verificação registradas.\\
Adicionado & Há comportamento implementado sem unidade correspondente no GDD completo.\\
Não verificável & Fonte, escopo ou evidência são insuficientes para decidir a correspondência.\\
\bottomrule
\multicolumn{2}{l}{\footnotesize Fonte: codebook operacional do projeto (Night Heist, 2026).}\\
\end{longtable}
\end{singlespace}

O instrumento distingue as ambiguidades lexical, conflitante, falta de parâmetro, escopo incerto e desconhecida, além de ausência de ambiguidade \rep. Para atribuição de responsável, admite discente, orientador, agente, conjunto ou desconhecido, sempre condicionado a registro original verificável \rep. As categorias são propostas para este caso e não constituem escala psicométrica ou instrumento previamente validado \rep.

\subsection{Procedimento de inspeção}
Primeiro, foram identificadas unidades explícitas e lacunas numéricas na leitura textual do GDD; depois, cada unidade foi associada a um critério observável antes da localização correspondente no código \fonte. A comparação conservou a diferença entre estrutura declarada e comportamento condicionado, de modo que uma constante foi usada como evidência estrutural e uma condição como evidência de regra lógica \rep.

Em seguida, foram examinados \texttt{scripts/level.gd}, \texttt{scripts/mission.gd} e \texttt{scripts/main.gd}, registrando símbolo ou método associado a cada unidade \rep. A classificação foi acompanhada de justificativa, sem deduzir que um trecho encontrado foi necessariamente executado em todas as plataformas previstas \rep. A última etapa do piloto consistiu em consolidar a matriz e seus limites; revisão independente e testes em runtime permanecem etapas futuras \rep.

\subsection{Ética, transparência e condições editoriais}
A unidade de análise é a decisão documentada relacionada ao código; a identificação de autoria foi fornecida pelo usuário, sem atribuição individual de desempenho ou de cada decisão \rep. O uso de produções escolares exige avaliar direitos e enquadramento pertinente, e este manuscrito não declara dispensa automática de avaliação ética apenas pela ausência de questionários (Brasil, 2016).

A concepção do recorte, a organização do instrumento, a inspeção documental e a redação deste manuscrito receberam assistência de ChatGPT/Codex; os resultados permanecem sujeitos à revisão intelectual dos pesquisadores humanos \rep. A autorização de continuidade do usuário suspendeu, nesta etapa de elaboração, a adaptação às condições da revista anteriormente considerada; não há declaração de conformidade editorial nem submissão realizada \rep. O repositório não possui licença geral do código no baseline auditado, portanto acesso público não foi apresentado como comprovação de licença open source do conjunto \rep.
